% =====================================================================
%  Multi-Reference EagleEye for Stellar Overdensity Detection in Gaia
%  Internal technical report
%  Compile:  pdflatex multireference_eagleeye_report.tex   (x2 for refs)
% =====================================================================
\documentclass[11pt,a4paper]{article}

\usepackage[T1]{fontenc}
\usepackage[utf8]{inputenc}
\usepackage{lmodern}
\usepackage{microtype}
\usepackage[margin=2.4cm]{geometry}
\usepackage{amsmath,amssymb}
\usepackage{booktabs}
\usepackage{graphicx}
\usepackage{xcolor}
\usepackage{caption}
\usepackage{array}
\usepackage{enumitem}
\usepackage{framed}
\usepackage[colorlinks=true,linkcolor=blue!50!black,citecolor=blue!50!black,
            urlcolor=blue!50!black]{hyperref}

\captionsetup{font=small,labelfont=bf}
% allow line breaks inside \texttt{} so long identifiers do not overflow
\renewcommand{\ttdefault}{lmtt}
\makeatletter\def\ttbreak{\hskip0pt plus0pt\relax}\makeatother
\emergencystretch=2em
\setlength{\parskip}{0.35em}
\setlength{\parindent}{0pt}

% --- shorthands -------------------------------------------------------
\newcommand{\unit}[1]{\ensuremath{\,\mathrm{#1}}}
\newcommand{\degsq}{\unit{deg^2}}
\newcommand{\masyr}{\unit{mas\,yr^{-1}}}
\newcommand{\pmvec}{\boldsymbol{\mu}}
\newcommand{\pmra}{\mu_{\alpha*}}
\newcommand{\pmdec}{\mu_{\delta}}
\newcommand{\Nchance}{\mathcal{N}_{\mathrm{ch}}}
\newcommand{\Yset}{\mathcal{Y}}
\newcommand{\Xset}{\mathcal{X}}

% ---- packages these rely on --------------------------------------------------
% \usepackage{amsmath,amssymb}
% \usepackage{bm}          % \bm gives proper bold math italic; \boldsymbol also works

% ---- vectors -----------------------------------------------------------------
\newcommand{\vb}[1]{\bm{#1}}                 % one place to change the vector style
\newcommand{\gvec}{\vb{g}}                   % density-gradient vector
\newcommand{\dvec}{\vb{d}}                   % reference-tile offset from Y's centre
\newcommand{\uvec}{\vb{u}}                   % position WITHIN a tile (local frame)
\newcommand{\xvec}{\vb{x}}                   % position in the field frame

% ---- sets --------------------------------------------------------------------
\newcommand{\Ycal}{\mathcal{Y}}              % test tile
\newcommand{\Xcal}{\mathcal{X}}              % reference tile; use \Xcal^{(j)}

% ---- statistics --------------------------------------------------------------
\newcommand{\Ups}{\Upsilon}                  % EagleEye per-reference score
\newcommand{\Gam}{\Gamma}                    % stacked score, \Gam = \sum_j \Ups^{(j)}
\newcommand{\eps}{\varepsilon}               % residual log density ratio
\newcommand{\Ocal}{\mathcal{O}}              % order symbol
\newcommand{\phat}{\hat{p}}                  % null Y-fraction n_Y/(n_Y+n_X)


\newcommand{\Ups}{\Upsilon}
\newcommand{\Gam}{\Gamma}
\newcommand{\Gstar}{\Gamma^{\star}}
\newcommand{\phat}{\hat{p}}
\newcommand{\Bhat}{\hat{B}}
\newcommand{\srb}{S/\sqrt{\Bhat}}
\newcommand{\Ycal}{\mathcal{Y}}
\newcommand{\Xcal}{\mathcal{X}}
\newcommand{\masyr}{\mathrm{mas\,yr^{-1}}}
\newcommand{\dvec}{\vec{d}}
\newcommand{\dg}{^{\circ}}
\newcommand{\Gaia}{\textit{Gaia}}

% --- callout box (framed + shaded, no mdframed/tcolorbox needed) -------
\definecolor{calloutbg}{gray}{0.955}
\definecolor{calloutrule}{gray}{0.70}
\newenvironment{callout}%
  {\par\smallskip
   \setlength{\fboxrule}{0.6pt}\setlength{\fboxsep}{7pt}%
   \def\FrameCommand{\fcolorbox{calloutrule}{calloutbg}}%
   \MakeFramed{\advance\hsize-\width \FrameRestore}}%
  {\endMakeFramed\par\smallskip}

\title{\bfseries Multi-Reference EagleEye for Stellar Overdensity\\
       Detection in \Gaia }
\author{}
\date{}

\begin{document}
\maketitle
\vspace{-3.2em}


\tableofcontents
\bigskip
\hrule

% =====================================================================
\section{Introduction}
\label{sec:purpose}

We apply EagleEye (Springer et al.\ 2025) to \Gaia\ DR3 phase-space
data to search for localised stellar overdensities.

Three things are established here. First, the \emph{multi-reference construction}: we compare a
test field against eight empirical neighbouring sky tiles rather than against a background model,
and combine the evidence into a single per-star statistic (Section~\ref{sec:construction}).
Second, \emph{reference handling}: a concentrated overdensity sitting in one of those reference
tiles corrupts the statistic for \emph{every} test star, and we describe the mechanism, its
detection, and its repair (Section~\ref{sec:refhandling}). Third, \emph{thresholding}: the
complete chain from raw \Gaia\ rows to a flagged cluster, with each threshold stated and justified
(Section~\ref{sec:thresholding}). We then show the pipeline recovering two known dwarf galaxies
(Section~\ref{sec:recovery}) and characterise one unidentified overdensity found alongside
Reticulum~II (Section~\ref{sec:anomaly}).

One thing is \textbf{not} established. We do not yet have a calibrated false-alarm rate, and
therefore we do not have a defensible detection threshold. Section~\ref{sec:calib} lists the
machinery that exists for this and the measurements that have not been made. Until those exist,
the per-cluster significance we report should be used for \emph{ranking} candidates, not for
deciding whether a candidate is real. This is stated in Section~\ref{sec:srb} where the statistic
is defined, and again in Section~\ref{sec:missing}.

Results in this report come from three runs, deliberately chosen to exercise different parts of
the pipeline. They were made at \textbf{three different configurations}, and are not directly
comparable with one another; Section~\ref{sec:config} gives the parameters and says where the
differences bite.

% =====================================================================
\section{EagleEye in brief}
\label{sec:ee}

This section is a summary of the parent method, included so that Section~\ref{sec:construction}
onwards can be read without reference to the EagleEye paper. Readers familiar with the method
should skip to Section~\ref{sec:inherited}.

\subsection{The coin-flip statistic}

Given a reference sample $X$ and a test sample $Y$, EagleEye scores each point by the composition
of its neighbourhood in the pooled set $U = X \cup Y$. For a test point $Y_i$, the membership of
its $k$-th nearest neighbour is encoded as $b_i^k = 1$ if that neighbour lies in $Y$ and $0$ if it
lies in $X$. Under the null hypothesis that the two samples are drawn from the same underlying
density, the cumulative count $B(i,K) = \sum_{k \le K} b_i^k$ behaves as a sequence of coin flips,
%
\begin{equation}
B(i,K) \sim \mathrm{Binomial}(K, \phat), \qquad \phat = \frac{n_Y}{n_X + n_Y}.
\end{equation}
%
The exact null is hypergeometric, because neighbours are drawn without replacement; the binomial
is an accurate approximation in the regime $K \ll n_X + n_Y$ that the method operates in.

The per-point anomaly score is the right-tail $p$-value of the observed count, maximised over
neighbourhood rank:
%
\begin{equation}
\Ups_i = \max_{1 \le K \le K_M}
         \left[ -\log \mathrm{pval}\!\left(B_{\rm obs}(i,K)\right) \right].
\label{eq:upsilon}
\end{equation}
%
Maximising over $K$ is what gives the statistic sensitivity across a range of anomaly scales
without committing to one. Points with $\Ups_i \ge \Ups^*_+$ are flagged, where the threshold is
calibrated by Monte Carlo (Section~\ref{sec:upsstar}).

The parameter $K_M$ is bounded above by the requirement that the neighbourhood be \emph{local}
--- that the $X\!:\!Y$ ratio be approximately constant across it. The paper sets
%
\begin{equation}
K_M \le 0.05 \min(n_X, n_Y).
\label{eq:locality}
\end{equation}
%
The $\min$ matters, and is re-derived operationally in Section~\ref{sec:km}; an expression of the
form $0.05(n_X + n_Y)$ appears elsewhere in the implementation but governs a different quantity,
the neighbour re-fetch buffer.

\subsection{IDE and rep\^echage}
\label{sec:ide}

The flagged set $Y_+$ contains a halo of points whose neighbourhoods merely \emph{overlap} a
genuine anomaly. Two refinement stages separate signal from halo.

\textbf{Iterative density equalisation (IDE)} repeatedly removes the highest-scoring remaining
point together with its nearest neighbours, and rescores, until nothing exceeds the threshold.
The removed set $\hat{Y}_+$ is a representation of the local density \emph{excess} --- the mass
that has to be taken out before the two samples become locally indistinguishable. It is smaller
than $Y_+$ and is the quantity used wherever we need ``the excess itself'' rather than
``everything the excess influenced''.

\textbf{Rep\^echage} clusters the flagged set using Density Peaks Advanced (d'Errico et al.\
2021), as implemented in \texttt{dadapy} (Glielmo et al.\ 2022), then, within each cluster
$\alpha$, sets a cluster-specific threshold at a low quantile ($q = 10^{-2}$) of the $\Ups$ values
of that cluster's IDE members. Points exceeding it form the final anomaly set
$Y^{\rm anom}_\alpha$.

The distinction matters later. Rep\^echage \textbf{deliberately re-admits the surrounding region}
and so contains background by construction. That is what makes the rep\^echage set usable as a
local background estimate (Section~\ref{sec:srb}) and it is why reference cleaning removes the
\emph{pruned} set and not the rep\^echaged one (Section~\ref{sec:equalise}).

\subsection{The injection background estimator}
\label{sec:eq9}

To estimate how much of an anomaly set is irreducible background, EagleEye moves each reference
point into the test set one at a time, rescores it, and counts how many are themselves flagged.
That injected count $Y^{\rm inj}_\alpha$, rescaled by the relative sizes of the two background
components, estimates the background in the anomaly region:
%
\begin{equation}
\Bhat_\alpha = \left|Y^{\rm inj}_\alpha\right|
               \frac{n_Y - |\hat{Y}_+|}{n_X - |\hat{X}_+|},
\qquad
\Lambda_\alpha = \frac{|Y^{\rm anom}_\alpha| - \Bhat_\alpha}{\sqrt{\Bhat_\alpha}}.
\label{eq:injection}
\end{equation}
%
$\Lambda_\alpha$ is interpretable as an approximate $z$-score only in the large-background regime;
outside it, it is a standardised measure of discrepancy strength. The estimator has a validity
boundary: when the pruned mass exceeds the opposing sample, the numerator of $\Bhat$ goes negative
and the statistic is undefined. This is rare in the balanced configurations used here
($\phat \approx 0.5$ throughout, Section~\ref{sec:recovery}) but common when references are pooled
(Section~\ref{sec:stackpool}).

\begin{callout}
\textbf{Placeholder.} A question concerning the implementation of this estimator is being raised
directly with the EagleEye authors. It is deliberately not documented here, and
Appendix~\ref{app:reserved} is reserved for it. Nothing in this report depends on its resolution:
all runs presented are balanced ($n_X \approx n_Y$), which is the regime least sensitive to it.
\end{callout}


\begin{table}[h]
\centering\small
\begin{tabular}{@{}p{0.46\textwidth}p{0.46\textwidth}@{}}
\toprule
\textbf{Inherited from EagleEye (Springer et al.\ 2025)} & \textbf{Introduced in this work}\\
\midrule
$\Ups$, the binomial $k$NN statistic & The eight-reference sky construction (\S\ref{sec:tiling})\\
IDE pruning; multimodal rep\^echage  & $\Gam = \sum_j \Ups^{(j)}$ and its reading as Fisher
                                       combination (\S\ref{sec:fisher})\\
The injection background estimator   & Reference equalisation (\S\ref{sec:equalise})\\
Monte Carlo calibration of $\Ups^*$  & The empirical $\Gstar$ bootstrap (\S\ref{sec:gstar})\\
                                     & $K_M$ resolution under the locality bound, including
                                       \texttt{"auto"} (\S\ref{sec:km})\\
                                     & Parallelisation with verified bit-level determinism
                                       (Appendix~\ref{app:parallel})\\
\bottomrule
\end{tabular}
\end{table}

One further item concerning the estimator implementation is being raised with the EagleEye authors
directly and is deliberately not documented here.

% =====================================================================
\section{Data, coverage and configuration}
\label{sec:data}

\subsection{The \Gaia\ query}

All data are from \texttt{gaiadr3.gaia\_source}. We retrieve \texttt{source\_id}, \texttt{ra},
\texttt{dec}, \texttt{pmra}, \texttt{pmdec}, \texttt{phot\_g\_mean\_mag}, \texttt{parallax},
\texttt{parallax\_error} and \texttt{ruwe}, adding \texttt{bp\_rp} when a colour--magnitude
diagram is required. The query is a square box in RA/Dec about the field centre, with RA wrapping
handled explicitly:

\begin{callout}
\ttfamily\footnotesize
SELECT source\_id, ra, dec, pmra, pmdec, phot\_g\_mean\_mag,\\
\hspace*{2em} parallax, parallax\_error, ruwe\\
FROM gaiadr3.gaia\_source\\
WHERE (ra BETWEEN \{ra\_min\} AND \{ra\_max\})\hfill -- or (ra >= ra\_min OR ra <= ra\_max)\\
\hspace*{2em} AND (dec BETWEEN \{dec\_min\} AND \{dec\_max\})\\
\hspace*{2em} AND pmra\hspace*{0.6em} BETWEEN -5 AND 5\\
\hspace*{2em} AND pmdec BETWEEN -5 AND 5\\
\hspace*{2em} AND (parallax - 3*parallax\_error) < 0.0
\end{callout}

Two cuts, each doing one job. The proper-motion box $|\mu| < 5\,\masyr$ removes the nearby disc
population, which would otherwise dominate the feature space, and bounds the PM features so that
their scatter is set by the halo rather than by a handful of high-velocity foreground stars. The
parallax cut $\varpi - 3\sigma_\varpi < 0$ retains only stars consistent with being distant, which
is the relevant population for satellites at tens of kiloparsecs.

Two things are deliberately \emph{not} cut, and both are open questions. We apply no RUWE cut, so
poorly-behaved astrometric solutions remain in the sample; and we apply no extinction correction,
so the stellar density field retains dust structure. The second is the more consequential, since a
dust-driven density gradient is exactly the kind of large-scale structure the reference
construction is meant to absorb (Section~\ref{sec:tiling}) and its residual is what
Sections~\ref{sec:refhandling} and~\ref{sec:scan} test for.

\subsection{Coverage: the precondition for every number in this report}
\label{sec:coverage}

The tiling routine does not complain when the $3\times3$ grid runs off the edge of the supplied
table. It simply returns thinner outer tiles. This is the most dangerous failure mode in the
pipeline, because a truncated reference tile contains too few stars, and test stars near the
corresponding edge are then compared against a region with almost no reference --- which the
statistic reads as an overwhelming overdensity pinned to the tile boundary.

This is not hypothetical. A run on a $\pm5\dg$ table with a $(-1,-1)$ grid offset clipped
\textbf{26 per cent of the western references' RA extent} and manufactured a 390-star ``cluster''
at $\srb = 239$, in a patch that raw star counts show to be slightly \emph{under}dense. The result
was invalidated entirely.

Two guards now exist. A hard coverage check inside the run driver raises if the supplied table does
not cover $\pm 3h$ about the window centre, where $h$ is the tile half-width. And the query can be
sized in advance: for a scan with box-size factors up to $s_{\max}$ and shift fraction $\rho$, the
half-width guaranteed to cover every configuration is
%
\begin{equation}
\mathrm{half} = h_0\left(1 + s_{\max}(3 + \rho)\right) + \mathrm{margin}.
\end{equation}
%
The operational rule, which we state as a rule because the incident above came from breaking it:
\textbf{never run the pipeline on a table you did not size in advance.}

\subsection{Feature construction}

Each tile is represented in coordinates local to its own centre,
%
\begin{equation}
\left[\Delta\alpha \cos\delta_0,\ \Delta\delta,\ \mu_{\alpha*},\ \mu_\delta\right],
\end{equation}
%
with $\Delta\alpha$ wrapped to $(-180\dg, 180\dg]$. The cosine factor uses the \emph{grid} centre
declination $\delta_0$ for all nine tiles rather than each tile's own declination. This is
deliberate: a per-tile cosine would give each tile a slightly different angular scale, and the
resulting feature spaces would not be commensurable when their $\Ups$ values are summed.

\subsection{Configuration}
\label{sec:config}

The three runs in this report use three different configurations. They are collected here as the
single source of truth; later sections refer to this table rather than restating parameters.

\begin{table}[h]
\centering\small
\caption{Configuration of the three runs presented. The runs are \emph{not} directly comparable.}
\begin{tabular}{@{}lccc@{}}
\toprule
Parameter & Bo\"otes~II (\S\ref{sec:bootes}) & Reticulum~II (\S\ref{sec:retii}, \S\ref{sec:anomaly})
          & Ret~II scan (\S\ref{sec:scan})\\
\midrule
Tile half-width $h$ [deg]      & 1.0000 & 1.3333 & 0.667--1.5\\
$K_M$                          & 50 & 100 & \texttt{"auto"} (32--169)\\
$p_{\rm ext}$                  & $10^{-3}$ & $5\times10^{-4}$ & $5\times10^{-3}$\\
$n_{\rm boot}$                 & 10 & 10 & 10\\
$N_{\rm null}$                 & 200\,000 & 200\,000 & 200\,000\\
DPA $Z$                        & 2.65 & 2.65 & 2.65\\
DBSCAN $\epsilon$ / min.\ samples & 1.0 / 5 & 1.0 / 5 & 1.0 / 5\\
Reference equalisation         & \textbf{on}, gate $\srb > 5$ & \textbf{off} & on\\
Seed                           & 42 & 42 & 42\\
\midrule
$n_Y$                          & 2\,099 & 3\,037 & 739--3\,548\\
$\sum_j n_X^{(j)}$             & 16\,517 & 23\,897 & ---\\
$\Gstar$                       & 43.79 & 52.06 & 35.2--45.1\\
Wall time                      & 32\,s & --- & ---\\
\bottomrule
\end{tabular}
\end{table}

\textbf{Two warnings follow from this table.}

\emph{The runs are not directly comparable.} $\Gstar$ differs between Bo\"otes~II and
Reticulum~II (43.79 against 52.06) principally because $p_{\rm ext}$ and $K_M$ differ, not because
the fields differ. Any comparison of significance between Sections~\ref{sec:retii}
and~\ref{sec:bootes} is therefore a comparison of two different tests, and we do not make one.

\emph{The Reticulum~II run has equalisation switched off.} This is a limitation of that run, not a
choice of method. We checked whether it matters: the nearest catalogued objects capable of
contaminating a reference tile --- NGC~1261, Horologium~I and Reticulum~III --- all fall
\textbf{outside} the $3\times3$ grid at $h = 1.3333\dg$ (NGC~1261 lies $5.85\dg$ west of the grid
centre in raw RA, against a grid half-extent of $4h = 4.0\dg$ in the tiling coordinate). There is
therefore no known contaminant in that field and no correctness problem. But it does mean the
Reticulum~II figures cannot illustrate equalisation, and the Bo\"otes~II run carries that load
instead (Sections~\ref{sec:refhandling} and~\ref{sec:bootes}).

% =====================================================================
\section{The multi-reference construction}
\label{sec:construction}

\subsection{The $3\times3$ tiling}
\label{sec:tiling}

We partition a square field into a $3\times3$ grid of tiles of half-width $h$. The central tile is
the test sample $\Ycal$; the eight surrounding tiles are reference samples $\Xcal^{(j)}$,
$j = 1,\dots,8$. Every reference is real sky, observed by the same instrument, subject to the same
selection function, at almost the same position on the sky. The fundamental assumption is that the local underlying background density is the same across all $\Xcal^{(j)}$. 

The arrangement is chosen so that the references bracket the test tile symmetrically. Four
antipodal pairs (E/W, N/S, and the two diagonals) surround $\Ycal$, so that a \emph{linear} density
gradient across the field contributes equally and oppositely to the members of each pair and
cancels in the mean. That cancellation is only partial, and it is important to say why immediately.
$\Gam = \sum_j \Ups^{(j)}$ is a sum of a non-negative, convex statistic. A linear gradient cancels
in the \emph{mean} of the underlying densities but not in the sum of the rectified scores: the tile on the sparser side contributes a positive $\Ups$ which the denser side cannot offset by
contributing a negative one, because $\Ups \ge 0$. A strong gradient therefore leaves a one-sided
residual along whichever edge of $\Ycal$ faces the sparser references. This is the motivation for
the geometric perturbation tests in Section~\ref{sec:scan} in which the test set tile $\Ycal$ is shifted in RA and DEC.


\subsubsection{Why antipodal bracketing cancels a gradient only partially}

\paragraph{Setup.}
Work in the local frame of the test tile, $\Ycal$ centred at the origin, with the
eight references $\Xcal^{(j)}$ at offsets $\dvec_j$ arranged in four antipodal pairs:
for each $j$ there is a $j'$ with $\dvec_{j'}=-\dvec_j$. Let the smooth surface
density over the $3\times3$ field be, to first order,
%
\begin{equation}
\rho(\xvec) \;=\; \rho_0\left(1+\gvec\cdot\xvec\right)+\Ocal(|\xvec|^2),
\qquad |\gvec\cdot\xvec|\ll1 .
\end{equation}
%
Each tile is normalised separately, since $\hat p_j=n_Y/(n_Y+n_{X_j})$ absorbs its
mean density. Writing $\uvec$ for position within a tile, the normalised densities are
%
\begin{equation}
f_Y(\uvec)=\frac{1+\gvec\cdot\uvec}{V},
\qquad
f_{X_j}(\uvec)=\frac{1}{V}\left(1+\frac{\gvec\cdot\uvec}{1+\gvec\cdot\dvec_j}\right),
\end{equation}
%
so the residual log-ratio that survives the per-tile normalisation is
%
\begin{equation}
\eps_j(\uvec)\;\equiv\;\log\frac{f_Y(\uvec)}{f_{X_j}(\uvec)}
\;=\;(\gvec\cdot\uvec)(\gvec\cdot\dvec_j)+\Ocal(g^3).
\label{eq:eps}
\end{equation}
%
The essential property is the sign structure, not the order: $\eps_j$ is
\emph{odd} in $\dvec_j$, hence
%
\begin{equation}
\eps_{j'}(\uvec)=-\eps_j(\uvec).
\label{eq:odd}
\end{equation}
%
(Without per-tile normalisation the residual would instead be $\Ocal(g)$, but still
odd in $\dvec_j$; everything below depends only on Eq.~\eqref{eq:odd}.)
Let $\Delta_j(\uvec)$ denote the corresponding standardised neighbour-composition
excess, so that $\Delta_{j'}=-\Delta_j$ likewise.

\paragraph{Lemma (the score is convex, not merely non-negative).}
For fixed $k$, $\Ups=-\log P\!\left(\mathrm{Bin}(k,\hat p)\ge n_k\right)$. Writing this
as $\psi(\Delta)=-\log Q(\Delta)$ with $Q$ the upper-tail function of the
null composition,
%
\begin{equation}
\psi'(\Delta)=\frac{q(\Delta)}{Q(\Delta)}\;=\;h(\Delta),
\end{equation}
%
the hazard function of the null. The binomial (and its Gaussian limit) is
log-concave, and the hazard of a log-concave density is non-decreasing; therefore
%
\begin{equation}
\psi\ge0,\qquad \psi'\ge0,\qquad \psi''\ge0 .
\label{eq:convex}
\end{equation}
%
Convexity is thus a theorem, not an assumption. Since $\Ups_i=\max_{k\in[20,K_M)}\psi_k$
and a pointwise maximum of convex functions is convex, Eq.~\eqref{eq:convex} survives
the scan over $k$.

\paragraph{Proposition (the pedestal).}
For an antipodal pair, Eqs.~\eqref{eq:odd} and \eqref{eq:convex} give
%
\begin{equation}
\Ups^{(j)}+\Ups^{(j')}=\psi(\Delta)+\psi(-\Delta)\;\ge\;2\,\psi(0),
\label{eq:jensen}
\end{equation}
%
by Jensen, with equality iff $\Delta=0$. The underlying densities cancel in the
\emph{mean} --- $\tfrac12(\bar\rho_j+\bar\rho_{j'})=\rho_0$ exactly at first order ---
but the sum of the \emph{rectified} scores does not. A linear statistic would give
$\Delta+(-\Delta)=0$; the convex one gives a strictly positive excess.

\paragraph{Two regimes.}
Expanding Eq.~\eqref{eq:jensen} about $\Delta=0$,
%
\begin{equation}
\Gam_{\mathrm{pair}}-2\psi(0)\;\simeq\;\psi''(0)\,\Delta^2+\Ocal(\Delta^4),
\label{eq:weak}
\end{equation}
%
so a \emph{weak} gradient contributes a pedestal that is \emph{quadratic} in the
gradient strength and \emph{symmetric}: both edges of $\Ycal$ along $\gvec$ are
affected equally. For a \emph{strong} gradient the asymmetry of $\psi$ takes over.
Since $\psi(\Delta)\sim\Delta^2/2$ as $\Delta\to+\infty$ while $\psi(\Delta)\to0$ as
$\Delta\to-\infty$,
%
\begin{equation}
\psi(\Delta)+\psi(-\Delta)\;\longrightarrow\;\psi(|\Delta|)
\qquad (|\Delta|\gtrsim1),
\end{equation}
%
i.e.\ the pair contributes \emph{only} through whichever reference makes $\Ycal$ look
overdense. By Eq.~\eqref{eq:eps} that is the sparser reference
($\gvec\cdot\dvec_j<0$), and its contribution is positive where
$\gvec\cdot\uvec<0$. The residual is therefore one-sided, concentrated on the edge of
$\Ycal$ facing the sparser references. Full rectification, not partial cancellation.

\paragraph{Corollary (stacking versus pooling).}
The same inequality separates the two combination schemes. Stacking rectifies then
sums; pooling averages the densities first and rectifies once. By Jensen,
%
\begin{equation}
\frac{1}{R}\sum_{j}\psi(\Delta_j)\;\ge\;\psi\!\left(\bar\Delta\right),
\qquad \bar\Delta=\frac1R\sum_j\Delta_j ,
\end{equation}
%
so a reference-level perturbation that averages to zero survives stacking and is
removed by pooling. This is the mechanism behind the measured contamination
asymmetry (stacked false-alarm inflation $75\times$ against pooled $37\times$): the
pedestal is additive and undiluted under $\Gam$.


\subsection{Robust scaling}
\label{sec:scaling}

Features are standardised by the median and median absolute deviation (MAD), fitted jointly across
all nine tiles. Fitting on the pooled set rather than on $\Ycal$ alone matters because the scaling
defines the metric in which nearest neighbours are found, and it must mean the same thing in the
test tile and in every reference.

The reason the scaling cannot simply be refitted per configuration is that it is not scale-free.
The positional MADs of a uniformly-filled tile scale linearly with the tile half-width, while the
proper-motion MADs are pinned near $1.4\,\masyr$ by the $\pm 5$ cut. Both runs confirm this
directly:

\begin{table}[h]
\centering\small
\caption{Measured MADs. The declination MAD is $0.50h$ in both runs; the RA MAD is
$0.50\,h\cos\delta_0$.}
\begin{tabular}{@{}lcccccccc@{}}
\toprule
Run & $h$ & MAD$(\Delta\alpha\cos\delta)$ & $/h$ & MAD$(\Delta\delta)$ & $/h$ &
MAD$(\mu_{\alpha*})$ & MAD$(\mu_\delta)$\\
\midrule
Bo\"otes~II   & 1.000 & 0.4894 & 0.489 & 0.4985 & 0.499 & 1.425 & 1.457\\
Reticulum~II  & 1.333 & 0.3943 & 0.296 & 0.6675 & 0.501 & 1.452 & 1.394\\
\bottomrule
\end{tabular}
\end{table}

The declination MAD is $0.50h$ in both cases, as expected for a uniform tile of half-width $h$.
The RA MAD is $0.50\,h\cos\delta_0$ --- $0.489 \approx 0.5\cos(12.9\dg)$ for Bo\"otes~II and
$0.296 \approx 0.5\cos(54.1\dg)$ for Reticulum~II. The positional MADs therefore vary by a factor
of $1.7$ between these two fields at fixed $h$ purely through declination, while the PM MADs are
constant to 4 per cent. The position-to-PM weighting of the Euclidean metric is consequently
\emph{not} a property of the method but of where and at what scale it is pointed.

Two consequences emerge. When comparing configurations at different $h$ --- as the robustness scan does
--- the metric must be frozen from one configuration and reused, or the scan measures a changing
distance function rather than a changing box. And because the DBSCAN radius $\epsilon$ is expressed
in these scaled units, freezing the metric is also what makes a fixed $\epsilon$ correspond to a
fixed angular size.

\subsection{Per-reference nulls}
\label{sec:perref}

Each reference has its own cardinality and therefore its own mixture proportion
$\phat_j = n_Y/(n_Y + n_X^{(j)})$. Each needs its own null distribution and its own threshold
$\Ups^*_j$. In the runs presented here the references are well balanced --- $\phat_j$ spans
$0.493$--$0.515$ for Bo\"otes~II and $0.490$--$0.519$ for Reticulum~II --- so the eight thresholds
are close, but they are computed separately regardless.

By contrast $K_M$ is shared across all eight comparisons, and must be. $\Gam$ sums the
$\Ups^{(j)}$, so they must be the same statistic; a per-reference $K_M$ would give each term a
different null \emph{and} a different attainable ceiling (Section~\ref{sec:ceiling}), so the sum
would be weighted by tile size rather than by evidence.

\subsection{$\Gam$ as Fisher combination}
\label{sec:fisher}

The combined statistic is
%
\begin{equation}
\Gam_i = \sum_{j=1}^{R} \Ups_i^{(j)}, \qquad R = 8.
\label{eq:gamma}
\end{equation}
%
Because $\Ups = -\log p$, this is $\Gam_i = -\log \prod_j p_i^{(j)}$, which is exactly Fisher's
combined-probability statistic (Fisher 1932). If the $R$ tests were independent with uniform
$p$-values, $2\Gam$ would be distributed as $\chi^2_{2R}$, equivalently
$\Gam \sim \mathrm{Gamma}(R,1)$.

\textbf{They are not independent.} All eight comparisons share the same test sample $\Ycal$, so the
Poisson noise of $\Ycal$ is common to all of them and the $\Ups^{(j)}$ are positively correlated by
construction. The nominal Fisher gain of $\sqrt{R}$ therefore overstates what stacking actually
buys. This is not a small concern: Guth \& Namjoo (2026) show that combining four correlated
anomaly tests of the CMB yields a joint $p$-value around $3\times10^{-8}$ that does not survive a
proper treatment of the dependence. The same arithmetic error is available to us here, and the
empirical null of Section~\ref{sec:gstar} is what we use instead of the analytic one precisely to
avoid it.

The degree of departure can be read off the fitted shape of the empirical $\Gam$ null. Matching a
$\mathrm{Gamma}(k,\theta)$ to its first two moments gives:

\begin{table}[h]
\centering\small
\begin{tabular}{@{}lccccc@{}}
\toprule
Run & mean & variance & $k_{\rm fit} = \mathrm{mean}^2/\mathrm{var}$ & $\theta$ & $R$\\
\midrule
Bo\"otes~II  & 13.38 & 52.76 & \textbf{3.39} & 3.94 & 8\\
Reticulum~II & 15.70 & 61.39 & \textbf{4.02} & 3.91 & 8\\
\bottomrule
\end{tabular}
\end{table}

so the null behaves like roughly four independent exponential contributions rather than eight.

\textbf{Two cautions on reading that number.} First, $k_{\rm fit}$ is measured on the
\emph{synthetic} uniform bootstrap null, so it characterises the method --- the shared-$\Ycal$
correlation and the max-over-$K$ selection --- and is close to the same number for any field. It
does not tell us whether \emph{our} references differ from one another. Second, $k_{\rm fit}$ is
not the same quantity as the effective reference count $R_{\rm eff} = R/[1 + (R-1)c]$ obtained from
the measured pairwise correlation $c$ of the $\Ups^{(j)}$ on real background stars. The two are
related by $\mathrm{mean}^2/\mathrm{var} = R_{\rm eff}\,(m^2/v)$ for marginal $\Ups$ mean $m$ and
variance $v$, and coincide only if $\Ups$ were exactly $\mathrm{Exp}(1)$, which it is not.

\textbf{$R_{\rm eff}$ is not reported here.} The per-reference $\Ups^{(j)}$ arrays are not retained
in the saved output of these runs, so the correlation cannot be recovered from them without
re-running. This is a gap in what we can currently claim --- it is the number that answers ``how
much independent evidence do eight references actually supply?'', and it is the honest counterpart
to the two-algorithm corroboration convention used elsewhere in the satellite-search literature
(e.g.\ Tan et al.\ 2026; Overdeck et al.\ 2026), which asserts agreement without quantifying it.
Retaining \texttt{Upsilon\_by\_ref} in the run output is a one-line change and is listed in
Section~\ref{sec:open}.

% =====================================================================
\section{Reference handling}
\label{sec:refhandling}

This section covers the failure mode that is specific to using real sky as a reference, and its
repair. It is illustrated throughout by the Bo\"otes~II field, where the contaminant is not a
hypothetical but a catalogued dwarf galaxy.

\subsection{Why a contaminated reference corrupts every test star}
\label{sec:pedestal}

$\Ups$ tests a \emph{local} neighbourhood composition against a \emph{global} mixture proportion
$\phat = n_Y/(n_Y + n_X)$. That comparison is valid as long as the two quantities correspond.

A reference tile that is uniformly denser than the test tile is handled correctly: the local
$X\!:\!Y$ ratio is elevated everywhere by the same factor that elevates $\phat$, and the two track
each other. But a \textbf{concentrated} clump --- a globular cluster, or, as here, a dwarf galaxy
--- is different. It inflates $n_X$, and therefore deflates $\phat$, while leaving the local
density at a randomly chosen field position completely unchanged. The global expectation and the
local reality decouple, and \emph{every} point of $\Ycal$ acquires an apparent excess of
test-sample neighbours relative to a $\phat$ that has been dragged down by structure nowhere near
it.

The effect is a pedestal added uniformly to $\Ups$ across the whole test tile. Measured on NGC~1261
in a Reticulum~II reference tile at a wider grid setting, $\phat$ fell from $0.524$ to $0.389$ and
a pedestal of $+4.66$ was added to $\Ups$ for all 7\,911 test stars. A pedestal of that size is
comparable to the detection threshold itself.

\subsection{The saturation ceiling}
\label{sec:ceiling}

There is a second, more insidious consequence. $\Ups$ is the negative log of a binomial right-tail
probability on a neighbourhood of at most $K_M$ points, so the largest value it can attain ---
realised when every usable neighbour belongs to the test set --- is
%
\begin{equation}
\Ups_{\max} = -(K_M - 1)\ln(1 - \phat).
\label{eq:ceiling}
\end{equation}
%
This ceiling depends only on $K_M$ and $\phat$. It is \textbf{independent of how large the
contaminant is}. As a reference grows relative to the test tile, $\phat \to 0$ and the ceiling
collapses.

When it collapses below the detection threshold $\Ups^*$, nothing in that reference can ever be
flagged, and the cleaning pass returns ``no clusters'' --- a result \textbf{indistinguishable from
a clean reference}. This has been observed: for M3 sitting in a Bo\"otes~III reference tile at
$K_M = 50$, $\Ups_{\max}$ and $\Ups^*$ were both $5.19$, the reference-side flagged set was empty,
and the most massive contaminant in the field was silently ignored.

The remedy is that the cleaning pass must run at its own $K_M$, chosen as large as the locality
bound permits, rather than at the $K_M$ chosen for the measurement. It must also refuse to proceed
rather than return a null result when even that leaves no headroom, and we require
$\Ups_{\max} - \Ups^* > 0.5$ before a pass is allowed to run.

\subsection{The equalisation algorithm}
\label{sec:equalise}

Equalisation cuts localised overdensities out of the reference tiles, expressed entirely in the
statistic's own terms. For each reference independently, and iterating:

\begin{enumerate}[nosep,leftmargin=1.4em]
\item Recompute $K_M$ from the current cardinalities as
      $\lfloor 0.05\min(n_X^{(j)}, n_Y)\rfloor$, subject to a floor of 25 (below which the usable
      rank range $20 \le K < K_M$ is empty).
\item Check the headroom $\Ups_{\max} - \Ups^*$ and raise if it is exhausted
      (Section~\ref{sec:ceiling}).
\item Run the full EagleEye pass with roles reversed, so that overdensities \emph{in the
      reference} are what is detected.
\item Gate the resulting reference-side clusters on $\srb > 5$.
\item Remove, for each gated cluster, its \textbf{IDE-pruned} set.
\end{enumerate}

Step 5 is the substantive choice. We remove the pruned set and not the rep\^echaged set, because
rep\^echage deliberately re-admits the surrounding region (Section~\ref{sec:ide}) and so contains
genuine background by construction. Removing only the excess mass excises the contaminant while
leaving the local background in place, so the repaired reference still samples the true background
at that position. Removing the rep\^echaged set instead would punch a hole in the reference.

There is deliberately \textbf{no cardinality target}. References of differing size are exactly what
the multi-reference construction is built to handle, and a tile that is legitimately denser should
remain denser. The loop stops when EagleEye no longer finds anything in the reference, not when the
tiles match.

\paragraph{A worked case.} In the Bo\"otes~II field, the reference tiles are not blank sky:
Bo\"otes~I, a kinematically confirmed dwarf with 204 members in the Battaglia et al.\ (2022)
catalogue, lies $1.66\dg$ north and $0.51\dg$ east of the grid centre. At $h = 1.0\dg$ this places
it in cell~7, the north-central \textbf{reference} tile --- outside $\Ycal$, but squarely inside
the background against which $\Ycal$ is measured. Only 6 of its 204 members reach the test tile.

Equalisation found it. The results for all eight references:

\begin{table}[h]
\centering\small
\caption{Reference equalisation in the Bo\"otes~II field. Seven references untouched; the one
containing Bo\"otes~I repaired.}
\begin{tabular}{@{}ccccccc@{}}
\toprule
Reference cell & $n$ before & $n$ after & cut & passes & $\srb$ & purity\\
\midrule
0 & 2\,079 & 2\,079 & 0 & 0 & --- & ---\\
1 & 2\,159 & 2\,159 & 0 & 0 & --- & ---\\
2 & 2\,152 & 2\,152 & 0 & 0 & --- & ---\\
3 & 2\,123 & 2\,123 & 0 & 0 & --- & ---\\
5 & 2\,043 & 2\,043 & 0 & 0 & --- & ---\\
6 & 1\,974 & 1\,974 & 0 & 0 & --- & ---\\
\textbf{7} & \textbf{2\,102} & \textbf{2\,005} & \textbf{97} & \textbf{1} & \textbf{26.02} &
\textbf{0.839}\\
8 & 1\,982 & 1\,982 & 0 & 0 & --- & ---\\
\bottomrule
\end{tabular}
\end{table}

Seven references were untouched; the one containing Bo\"otes~I was flagged at reference-side
$\srb = 26.0$ and 97 stars were excised in a single pass, at an estimated purity of 0.84. The
repaired tile is drawn in Figure~\ref{fig:bootes}, where the excised stars are shown separately.

This is the clearest available demonstration that the failure mode is real, that it arises in the
ordinary course of pointing at a field with known neighbours rather than in a contrived case, and
that the method detects its own contamination without being told where to look.

\subsection{$K_M$ and the locality bound}
\label{sec:km}

$\Ups$ assumes the $X\!:\!Y$ ratio is constant inside the $K$-neighbourhood. It is not --- two
adjacent sky tiles have a slowly varying ratio --- and the neighbourhood radius grows like
$k^{1/d}$, so the larger $K_M$ is, the more of that variation the statistic swallows and reports as
signal. On a synthetic field with a mild ratio gradient and no injected anomaly, the false-positive
rate as a multiple of the nominal $p_{\rm ext}$ is:

\begin{table}[h]
\centering\small
\caption{False-positive inflation as a multiple of nominal $p_{\rm ext}$, measured on a synthetic
gradient field with no anomaly.}
\begin{tabular}{@{}lccccc@{}}
\toprule
$K_M$ & 25 & 50 & 100 & 300 & 600\\
\midrule
$d = 2$ & 2.0$\times$ & 6.7$\times$ & 12.5$\times$ & 46.0$\times$ & 118$\times$\\
$d = 4$ & 1.3$\times$ & 1.8$\times$ & 2.0$\times$ & 2.2$\times$ & 6.8$\times$\\
\bottomrule
\end{tabular}
\end{table}

The features here are four-dimensional, where the inflation is far milder --- a factor of two out
to $K_M = 300$. This is why setting $K_M$ to the locality ceiling is a reasonable default rather
than a reckless one.

Pushing against that, larger $K_M$ is broadly \emph{better} for sensitivity: the threshold $\Ups^*$
grows only logarithmically with $K_M$ while the ceiling $\Ups_{\max}$ grows linearly, and a diffuse
anomaly accumulates its excess over many neighbours and may be invisible at small $K_M$ regardless
of how many stars it contains. The \texttt{"auto"} setting resolves $K_M$ to the largest value the
locality bound permits, and it is resolved \emph{after} equalisation --- resolving it before would
let a contaminated tile's inflated cardinality set the bound.

The consequences of the bound are visible in the robustness scan of Section~\ref{sec:scan}, where
shrinking the box by a factor $0.67$ reduces $n_Y$ to 739 and drives $K_M^{\rm auto}$ down to 32,
close to the floor of 25.

% =====================================================================
\section{Thresholding: the complete chain}
\label{sec:thresholding}

This section states every threshold between raw \Gaia\ rows and a flagged cluster. The chain is
%
\begin{equation*}
\text{rows} \rightarrow \Ycal,\ \Xcal^{(j)} \rightarrow \Ups^*_j \rightarrow \Ups_i^{(j)}
\rightarrow \Gam_i \rightarrow \Gstar \rightarrow A_\Gam \rightarrow \text{DBSCAN}
\rightarrow \Bhat_{\rm w} \rightarrow \srb.
\end{equation*}

\subsection{$\Ups^*_j$ --- the per-reference flagging threshold}
\label{sec:upsstar}

Because $\Ups_i$ is a maximum over $K$, its null distribution is not the tail of any single
binomial and must be obtained by simulation. For each reference we draw $N = 200\,000$ independent
Bernoulli sequences of length $K_M$ at probability $\phat_j$, compute $\Ups$ for each exactly as for
a real point, and take
%
\begin{equation}
\Ups^*_j = \mathrm{quantile}\left(\{\Ups\},\ 1 - p_{\rm ext}\right).
\end{equation}
%
One threshold per reference, because $\phat_j$ differs.

\textbf{$p_{\rm ext}$ is a per-point exceedance level.} It is not a cluster-level or field-level
error rate, and it does not describe how often the pipeline produces a spurious \emph{cluster}.
This is the most commonly misread parameter in the pipeline and it is the reason
Section~\ref{sec:calib} exists: converting a per-point level into a statement about candidate
reliability requires the false-alarm measurements we have not yet made.

\subsection{$\Gstar$ --- the empirical threshold}
\label{sec:gstar}

The threshold on the combined statistic is obtained by bootstrap. For each of
$n_{\rm boot} = 10$ realisations we draw $Y_b$ and $X_b^{(j)}$ uniformly on $[0,1]^4$ \textbf{at the
realised cardinalities of the actual run}, push them through the identical multi-reference
pipeline, and collect the resulting $\Gam_i^{(b)}$. Pooling across realisations and points gives the
null sample, and
%
\begin{equation}
\Gstar = \mathrm{quantile}\left(\Gam_{\rm null},\ 1 - p_{\rm ext}\right).
\end{equation}

For the two runs here this gives null samples of $20\,990$ and $30\,370$ values respectively, and:

\begin{table}[h]
\centering\small
\caption{Empirical $\Gam$ null quantiles.}
\begin{tabular}{@{}lcccccc@{}}
\toprule
Run & $q_{0.9}$ & $q_{0.99}$ & $q_{0.999}$ & $q_{0.9999}$ & $\Gstar$ used & $\max_i \Gam_i$\\
\midrule
Bo\"otes~II & 23.32 & 35.06 & 43.79 & 49.28 & \textbf{43.79} ($p_{\rm ext}=10^{-3}$) & 76.22\\
Reticulum~II & 26.24 & 38.80 & 49.56 & 57.46 & \textbf{52.06} ($p_{\rm ext}=5\times10^{-4}$) &
253.08\\
\bottomrule
\end{tabular}
\end{table}

The empirical route is used rather than the analytic $\mathrm{Gamma}(k_{\rm fit},\theta)$ fit
because the true law is a sum of correlated maxima-of-exponentials whose tail the max-over-$K$
selection reshapes. The fitted form tracks the null well to about the $10^{-2}$ level and then runs
high --- by roughly 9 per cent at $p_{\rm ext} = 10^{-3}$ and 15 per cent at $10^{-4}$.
Substituting it would buy better run-to-run stability at the cost of a larger bias, in the
conservative direction, discarding real anomalies.

The tail is where $n_{\rm boot}$ matters. At $n_{\rm boot} = 10$ and
$n_Y \approx 2\,000$--$3\,000$ the null sample contains $2$--$3\times10^4$ values, so a quantile at
$p_{\rm ext} = 10^{-3}$ is estimated from roughly 20--30 points and at $5\times10^{-4}$ from roughly
10--15. These are thin, and $\Gstar$ carries a corresponding sampling uncertainty that we have not
propagated. Running at smaller $p_{\rm ext}$ without increasing $n_{\rm boot}$ is not meaningful.

\subsection{$A_\Gam$ and clustering}

Points exceeding the threshold form the anomalous set
$A_\Gam = \{i : \Gam_i > \Gstar\}$, which is then partitioned by DBSCAN (Ester et al.\ 1996) in the
full four-dimensional scaled feature space with $\epsilon = 1.0$ and min.\ samples $= 5$.
Clustering in the same space in which the statistic was computed --- rather than on the sky alone
--- means a group must be coherent in position \emph{and} proper motion to survive.

Two properties of this step should be kept in mind. Because $\epsilon$ is in MAD units, its
physical meaning is fixed only if the metric is frozen (Section~\ref{sec:scaling}). And DBSCAN's
noise label is retained rather than suppressed: points above threshold that join no cluster are
reported as unclustered and excluded from the cluster statistics. For Bo\"otes~II, 20 points
exceeded $\Gstar$ and 18 formed a single cluster, leaving 2 unclustered; for Reticulum~II, 111
points exceeded threshold and all 111 were assigned to two clusters with no unclustered remainder.

\subsection{$\Bhat$ and $\srb$ per cluster}
\label{sec:srb}

Each $\Gam$ cluster needs a background estimate, and the estimate available is per-reference: the
rep\^echage stage of each of the eight comparisons produces its own clusters, each with its own
$\Bhat$ from the injection estimator of Section~\ref{sec:eq9}. We combine them by overlap
weighting. For a $\Gam$ cluster $\alpha$, let $o^{(j)}_\alpha$ be the number of stars it shares with
the overlapping rep\^echage cluster in reference $j$. Then
%
\begin{equation}
\Bhat_\alpha = \frac{\sum_j o^{(j)}_\alpha \Bhat^{(j)}_\alpha}{\sum_j o^{(j)}_\alpha},
\qquad S = |A_\Gam^{(\alpha)}|,
\qquad \text{reported as } S/\sqrt{\Bhat_\alpha}.
\end{equation}
%
Only references contributing a finite, positive $\Bhat$ and a finite, positive significance are
included, so the number of contributing references is itself diagnostic: a cluster supported by all
eight is a different object from one supported by two. In Section~\ref{sec:recovery} the
Reticulum~II detection draws on all 8 references while the unidentified cluster draws on 2, and the
Bo\"otes~II detection on 5.

\textbf{What this number is, and is not.} It is a local estimate of excess over background, and it
is a useful \emph{ranking} statistic. It is not a calibrated significance. It takes no account of
how many independent places in the tile could have produced a comparable clump by chance, so it
cannot be read as a trials-corrected detection significance, and it should not be quoted as a
number of sigma. Converting it into such a statement is what Section~\ref{sec:calib} is for, and
that work is not done.

% =====================================================================
\section{Recovery of known systems}
\label{sec:recovery}

We show the pipeline recovering two kinematically confirmed dwarf galaxies from the
Drlica-Wagner et al.\ (2020) catalogue, using membership from Battaglia et al.\ (2022) at
$P_{\rm memb} > 0.5$ as ground truth. Members are matched to our \Gaia\ DR3 rows by
\texttt{source\_id}, so the comparison involves no positional tolerance.

For context, an earlier batch application of this pipeline to all 30 class-4 dwarfs with usable
fields recovered 21 at a true-positive rate above 0.5 and 9 not at all, with no intermediate
outcomes. That batch used a different configuration from anything in this report (wider tiles,
$K_M = 300$, $p_{\rm ext} = 10^{-5}$, no equalisation) and its per-system numbers are therefore not
comparable with those below; Bo\"otes~II in particular is re-run here under the configuration of
Section~\ref{sec:config} and the result differs. We note the aggregate only as an indication of
scale and do not draw on the batch further.

\subsection{Reticulum~II}
\label{sec:retii}

Reticulum~II has 75 Battaglia members, all of which fall inside the test tile. The run returns
\textbf{two} $\Gam$ clusters from 111 points above threshold, with no unclustered remainder.

\begin{table}[h]
\centering\small
\caption{The two $\Gam$ clusters in the Reticulum~II field.}
\begin{tabular}{@{}lcc@{}}
\toprule
 & cluster 1 & cluster 0\\
\midrule
$n$                                   & 89 & 22\\
$\Bhat$ (overlap-weighted)            & 19.96 & 5.34\\
$\srb$                                & 19.92 & 9.52\\
Contributing references               & 8 of 8 & 2 of 8\\
Best catalogue match                  & \textbf{Reticulum~II} & \emph{unidentified}\\
Members recovered                     & 62 of 75 & 0\\
TPR                                   & \textbf{0.827} & ---\\
Purity against known members          & 0.697 & 0.000\\
Centroid $(\alpha,\delta)$            & $(53.908, -54.036)$ & $(52.884, -55.027)$\\
Separation from catalogue position    & $0.017\dg$ & $1.142\dg$\\
Mean proper motion [$\masyr$]         & $(+2.42, -1.44)$ & $(-0.40, -0.01)$\\
PM dispersion [$\masyr$]              & $(0.38, 0.44)$ & $(0.38, 0.35)$\\
Median $G$                            & 19.68 & 20.29\\
\bottomrule
\end{tabular}
\end{table}

Cluster~1 is Reticulum~II. It recovers 62 of 75 known members, its centroid lies $1.0'$ from the
catalogue position, and its mean proper motion matches the published value. It is supported by all
eight references, which is what a genuine, strong overdensity should look like: no choice of
background makes it go away. Of its 89 members, 27 are not in the Battaglia list --- a purity of
0.70 --- which is the expected behaviour of the rep\^echage stage, which admits the surrounding
region by construction (Section~\ref{sec:ide}), and is not by itself evidence of contamination.

Cluster~0 matches nothing in the catalogue and is the subject of Section~\ref{sec:anomaly}.

\begin{figure}[htbp]
\centering
\includegraphics[width=0.78\textwidth]{figures/fig1_reticulumII_sky.png}
\caption{The Reticulum~II test tile, with the 111 points above $\Gstar$ marked and the two clusters
distinguished. Reticulum~II's Battaglia members are overplotted as red stars. Cluster~1 sits on the
target at the tile centre; cluster~0 occupies the south-western corner. Proper-motion counterpart:
\texttt{figures/fig1b\_reticulumII\_pm.png}.}
\label{fig:retii}
\end{figure}

\subsection{Bo\"otes~II}
\label{sec:bootes}

Bo\"otes~II has 22 Battaglia members, all inside the test tile. The run returns a single $\Gam$
cluster from 20 points above threshold, with 2 unclustered.

\begin{table}[h]
\centering\small
\caption{The single $\Gam$ cluster in the Bo\"otes~II field.}
\begin{tabular}{@{}lc@{}}
\toprule
 & value\\
\midrule
$n$                                & 18\\
$\Bhat$ (overlap-weighted)         & 3.93\\
$\srb$                             & 9.08\\
Contributing references            & 5 of 8\\
Best catalogue match               & \textbf{Bo\"otes~II}\\
Members recovered                  & 14 of 22\\
TPR                                & \textbf{0.636}\\
Purity against known members       & 0.778\\
Centroid $(\alpha,\delta)$         & $(209.537, +12.861)$\\
Separation from catalogue position & $0.023\dg$\\
Mean proper motion [$\masyr$]      & $(-2.54, -0.36)$\\
\bottomrule
\end{tabular}
\end{table}

Bo\"otes~II is a considerably harder target than Reticulum~II --- 22 catalogued members against 75
--- and the recovery is correspondingly less complete: 14 of 22 members, at $\srb = 9.08$,
supported by 5 of the 8 references rather than all 8. The centroid is $1.4'$ from the catalogue
position. Of the 18 cluster members, 14 are known members, so the cluster is 78 per cent pure.

The more instructive aspect of this field is what is happening in the \emph{references}. Bo\"otes~I
lies $1.66\dg$ north of the grid centre, placing it inside the north-central reference tile
(Section~\ref{sec:equalise}), with 204 catalogued members of which only 6 reach the test tile.
Equalisation flagged that tile at reference-side $\srb = 26.0$ and removed 97 stars; the other
seven references were untouched. Without that repair, Bo\"otes~I's 204 members would have inflated
$n_X$ for one of the eight comparisons, depressing $\phat$ and adding a pedestal to $\Ups^{(j)}$ for
all 2\,099 test stars (Section~\ref{sec:pedestal}).

This is the case the multi-reference construction has to handle in order to be usable: any field
chosen because it contains one known satellite is reasonably likely to contain another within a
few degrees.

\begin{figure}[htbp]
\centering
\includegraphics[width=0.74\textwidth]{figures/fig2_bootesII_sky.png}
\caption{The Bo\"otes~II test tile with the recovered cluster, the Battaglia members of Bo\"otes~II
(targets) and of Bo\"otes~I (interlopers, 6 inside the tile and 198 outside), and the repaired
reference tile drawn at the top with its 97 excised stars marked. Colour--magnitude counterpart:
\texttt{figures/fig2b\_bootesII\_cmd.png}.}
\label{fig:bootes}
\end{figure}

% =====================================================================
\section{An unidentified overdensity in the Reticulum~II field}
\label{sec:anomaly}

Cluster~0 of Section~\ref{sec:retii} matches no catalogued system. This section characterises it.
We present it as a \textbf{candidate requiring follow-up, not as a detection}, and the evidence
below is mixed.

\subsection{Properties}

The cluster contains 22 stars at $(\alpha, \delta) = (52.884, -55.027)$, $1.14\dg$ from
Reticulum~II. It is compact and round --- its members span 0.26 and 0.28 of the tile in RA and
declination respectively, an aspect ratio of 1.1:1, which is essentially identical to
Reticulum~II's own footprint in the same run ($0.27 \times 0.29$, 1.1:1).

Its proper motion is tightly concentrated at $(-0.40, -0.01)\,\masyr$ with dispersions
$(0.38, 0.35)$, again comparable to Reticulum~II's $(0.38, 0.44)$. Its members are about 0.6\,mag
fainter in median $G$ (20.29 against 19.68).

\textbf{It is not associated with Reticulum~II.} The two proper motions differ by
$2.9\,\masyr$, an order of magnitude larger than either dispersion. At Reticulum~II's distance of
roughly 30\,kpc that corresponds to a relative transverse velocity of several hundred
km\,s$^{-1}$, far in excess of anything bound to it.

An overdensity at a consistent position has been seen in this field at other configurations, with
somewhat different membership --- 27 stars at $(52.881, -54.936)$ in an earlier run. The right
ascensions agree to $0.003\dg$ and the declinations differ by $0.09\dg$. We take these to be the
same feature seen under different tilings, but we have not established that by matching
\texttt{source\_id} between runs, and it should be checked.

\subsection{Where the evidence is weak}

Two things should be stated plainly before the robustness results.

\textbf{It is only supported by 2 of the 8 references}, against 8 of 8 for Reticulum~II in the same
run. The overlap-weighted background is therefore built from two rep\^echage clusters rather than
eight, and $\Bhat = 5.34$ is correspondingly less well determined than Reticulum~II's 19.96.

\textbf{It touches the tile boundary.} Its members come within $0.011h$ of the western edge and
$0.001h$ of the southern edge --- that is, they reach the corner of the test tile. This matters in
two ways. Its true extent may be larger than measured, because any members beyond the boundary are
not in $\Ycal$ at all; and worse, those members would be sitting in the adjacent reference tiles,
where they would act to suppress the very feature we are trying to measure. The 22 stars should be
read as a lower bound on the membership.

It is \emph{not}, however, the classic tiling artefact. That signature is a thin strip lying flush
along an edge and spanning the full width of the tile --- high aspect ratio, large span, centroid
pushed to the boundary. This cluster has an aspect ratio of 1.1 and spans about a quarter of the
tile in each direction. It is a compact clump that happens to lie near a corner, which is a
different thing, and the standard geometric test does not flag it.

\subsection{The robustness scan}
\label{sec:scan}

To separate a feature locked to the \emph{sky} from one locked to the \emph{tiling}, the pipeline
is re-run on eight grids anchored on the anomaly itself: the nominal grid, four shifts pushing the
anomaly toward each tile edge, two changes of box size, and one diagonal shift. A real object stays
at a fixed sky position while the grid moves under it; an artefact follows the grid.

The reference set $A_0$ for this scan contains 28 stars.

\begin{table}[h]
\centering\footnotesize
\caption{The eight-configuration robustness scan. Six of seven perturbed configurations recover the
feature; the failure is the smallest box, where the locality bound drives $K_M$ to 32.}
\begin{tabular}{@{}lrrrrrcrrrr@{}}
\toprule
Configuration & $h$ & $n_Y$ & $K_M$ & $\Gstar$ & clusters & recovered & overlap & recovery &
Jaccard & $\srb$\\
\midrule
nominal            & 1.000 & 1\,627 & 75 & 39.75 & 3 & yes & 27 & 0.96 & 0.93 & 8.63\\
shiftE             & 1.000 & 1\,659 & 75 & 40.43 & 2 & yes & 22 & 0.79 & 0.79 & 6.76\\
shiftW             & 1.000 & 1\,585 & 74 & 40.53 & 1 & yes & 26 & 0.93 & 0.84 & 6.53\\
shiftN             & 1.000 & 1\,577 & 74 & 40.65 & 2 & yes & 24 & 0.86 & 0.80 & 6.63\\
shiftS             & 1.000 & 1\,627 & 75 & 39.94 & 2 & yes & 27 & 0.96 & 0.79 & 7.78\\
\textbf{scale $\times$0.67} & 0.667 & \textbf{739} & \textbf{32} & 35.22 & \textbf{0} &
\textbf{no} & 0 & 0.00 & 0.00 & ---\\
scale $\times$1.50 & 1.500 & 3\,548 & 169 & 45.09 & 3 & yes & 23 & 0.82 & 0.64 & 6.98\\
diagonal           & 1.000 & 1\,596 & 74 & 39.01 & 2 & yes & 27 & 0.96 & 0.84 & 9.00\\
\bottomrule
\end{tabular}
\end{table}

\textbf{The formal verdict is ROBUST.} Six of the seven perturbed configurations recover the
feature (0.86, against a requirement of 0.75); the median Jaccard index among recoveries is 0.80,
against a requirement of 0.5; and the significance ranges over 6.53--9.00, a ratio of 1.38, against
a permitted ratio of 3.0. The feature is locked to the sky, not to the tiling.

The single failure is the smallest box, and it is explained by the method rather than being
anomalous. Shrinking the tile to $h = 0.667\dg$ reduces $n_Y$ to 739, which through the locality
bound of Section~\ref{sec:km} drives $K_M^{\rm auto}$ down to 32 --- close to the floor of 25 below
which no test exists at all. The attainable $\Ups_{\max}$ falls with it, and a feature of this
strength can no longer clear the threshold: 18 per cent of $A_0$ remained above $\Gstar$ but no
cluster formed. This is the documented behaviour of a shrinking box, not evidence against the
feature. It does mean that a single significance threshold applied across all eight configurations
would be biased against the small ones.

\begin{figure}[htbp]
\centering
\includegraphics[width=\textwidth]{figures/fig3_scan_overlay.png}
\caption{Left: the sky, with the recovered anomaly's contour drawn for each configuration and each
configuration's test-tile boundary shown dashed. The contours coincide while the boundaries move,
which is the signature of a sky-locked feature. Right: the same in proper-motion space, where
Reticulum~II is visible as the field overdensity near $(+2.4, -1.4)\,\masyr$ and the anomaly
contours sit separately near $(-0.4, 0.0)$. The $\times 0.67$ configuration is absent from both
panels because it returned no cluster.}
\label{fig:scan}
\end{figure}

\subsection{What has been ruled out, and what has not}

Two interpretations have been assessed and set aside. A stellar-wake origin --- an overdensity of
halo field stars trailing a dark perturber --- is not supported: the shear from any plausible
perturber at this distance is orders of magnitude below what is measurable with these data, and the
proper motions are in any case incompatible with a common origin with Reticulum~II. A boundedness
argument is also unavailable: the test has no discriminating power at 30\,kpc, as demonstrated by
running it on Reticulum~II itself, a known bound system, where it fails identically. Both are
reported so they are not attempted again.

What has \textbf{not} been established is the alternative that matters: that 22 stars this tightly
grouped in proper motion could not arise by chance within a field whose proper-motion distribution
is itself strongly clustered. Answering that requires the calibration of Section~\ref{sec:calib},
and until it exists the candidate stands on the robustness scan alone --- which shows the feature
is real in the data and not an artefact of the tiling, but not that it is a physical system.

% =====================================================================
\section{Calibration --- PARTIAL, FOR LATER COMPLETION}
\label{sec:calib}

\begin{callout}
This section is deliberately incomplete. It records the machinery that exists and the measurements
that have not been made. \textbf{No threshold recommendation should be taken from it.}
\end{callout}

\subsection{Nulls implemented}
\label{sec:nulls}

Four nulls exist, in increasing order of realism. Each answers a different question and none
answers all of them.

\textbf{Proper-motion scramble.} Permute the $(\mu_{\alpha*}, \mu_\delta)$ pairs among stars in the
window. Positions are untouched, so the spatial density field and all its gradients survive; the
joint PM distribution survives because the pairs are only relabelled; what is destroyed is the
position--PM coherence that a genuine comoving group has. Applied to the Reticulum~II-field
candidate this gave an observed $\srb = 7.05$ against a maximum of 3.54 over 12 blank trials,
$p < 0.08$. Twelve trials is too few to say more. This null can only answer ``could the pipeline
have manufactured this from the marginals alone?'' --- it cannot distinguish a genuine chance
density from a physical system.

\textbf{Off-position.} Run the identical pipeline at positions elsewhere on the sky and read the
empirical distribution of $\srb$ off blank sky. This manipulates nothing and is the classic
off-source measurement. Two hardening steps were necessary after a first pass was found unsound:
the exclusion radius around known objects must be the full $3h$ half-width of the grid rather than
the test tile, because an object just outside $\Ycal$ sits in a \emph{reference} and does exactly
the damage described in Section~\ref{sec:pedestal}; and the DW20 catalogue contains no globular
clusters, so those must be excluded using a separate catalogue (Vasiliev \& Baumgardt 2021).
Off-positions are also matched in stellar density to the target field, since surface density varies
strongly with Galactic latitude.

\textbf{Smooth mock.} Resample the field with replacement and jitter by a kernel wider than any
clump but narrower than the field's real structure. Below the bandwidth the generating density is
smooth by construction, so anything recovered is manufactured by the method, while the large-scale
gradient and tile-to-tile density differences survive. This isolates method artefacts --- tile-edge
effects, the rectification described in Section~\ref{sec:tiling}, DBSCAN scale coupling --- without
appealing to what is or is not in the sky.

\textbf{Scan-level persistence.} The one that matters for flagging policy. A marginal cluster in a
single configuration means little; the same cluster reappearing at the same sky position across all
eight configurations of Section~\ref{sec:scan} is a different claim, and carries power the
single-run significance does not express. Measuring how often blank fields produce such a
persistent group gives the number to flag on. The machinery is written and parallelised.
\textbf{No number has been produced.}

\subsection{Stacking against pooling}
\label{sec:stackpool}

There is an alternative use of eight references: concatenate them into a single pooled reference
and run the standard single-comparison pipeline once. This is a genuinely different statistic, with
$\phat \approx 1/9$ rather than $\approx 1/2$, and it is worth recording as a limitation of
stacking that the two are not uniformly ordered.

On synthetic fields built with the real cardinalities of a \Gaia\ tile and deliberately lumpy,
unequal-sized references, recovery rates were:

\begin{table}[h]
\centering\small
\begin{tabular}{@{}lcc@{}}
\toprule
Anomaly & Stacked & Pooled\\
\midrule
Compact ($w = 0.035$, $N = 10$) & 0.25 & \textbf{0.95}\\
Diffuse ($w = 0.085$, $N = 20$) & \textbf{0.43} & 0.23\\
\bottomrule
\end{tabular}
\end{table}

Compact anomalies favour pooling and diffuse anomalies favour stacking. The mechanism is
understandable: stacking is a persistence measure and suppresses variance on a signal present in
every comparison, which is what a diffuse excess looks like; but a $\phat$ pedestal arising from
reference lumpiness is additive and undiluted under stacking, whereas pooling averages it away.

\textbf{These numbers rest on two trials per grid cell and should not be relied on.} The direction
of the effect is suggestive and the mechanism is plausible, but the statistics are thin. We record
it here as a caveat attached to the stacked approach used throughout this report, and note that the
candidate of Section~\ref{sec:anomaly} --- compact, 22 stars --- sits in the regime where stacking
performed worse in this test. Running the pooled comparison on the real field is a cheap
cross-check and has not been done.

\subsection{What is missing}
\label{sec:missing}

\begin{itemize}[nosep,leftmargin=1.4em]
\item The scan-level persistence false-alarm rate. Staged and parallelised; \textbf{no number}.
\item The off-position null re-run under the hardened exclusion logic.
\item The scramble null beyond 12 trials.
\item Whether the stacking/pooling crossover survives more than two trials per cell.
\item A pooled cross-check on the Reticulum~II field.
\end{itemize}

\textbf{Consequence.} Until these exist there is no defensible detection threshold, and $\srb$
should be used to rank candidates, not to decide them. The candidate of Section~\ref{sec:anomaly}
is reported on that basis.

% =====================================================================
\section{Open items}
\label{sec:open}

\paragraph{Code.}
\begin{itemize}[nosep,leftmargin=1.4em]
\item The per-reference $\Ups^{(j)}$ arrays are not retained in the saved run output, so
      $R_{\rm eff}$ cannot be computed after the fact (Section~\ref{sec:fisher}). Retaining
      \texttt{Upsilon\_by\_ref} is a one-line change and should be made before the next production
      runs.
\item The injection estimator divides by zero rather than returning NaN when the injected count is
      zero.
\item One further item concerning that estimator is pending discussion with the EagleEye authors
      (Sections~\ref{sec:eq9} and Appendix~\ref{app:reserved}).
\item The scan-overlay plotting has unguarded RA-wrap exposure.
\item The null-calibration routine is unseeded.
\item \textbf{Nothing is committed.} The EagleEye working copy carries uncommitted changes to
      \texttt{EagleEye.py}, \texttt{utils\_EE.py}, \texttt{utils\_sistematics.py} and
      \texttt{test\_multiBkg\_ordered.py}. This should be resolved before anything else touches
      those files, since the runs in this report depend on that working state and it is not
      currently reproducible from version control.
\end{itemize}

\paragraph{Analysis.}
\begin{itemize}[nosep,leftmargin=1.4em]
\item Confirm by \texttt{source\_id} that the 22-star cluster of Section~\ref{sec:anomaly} and the
      27-star cluster seen at other configurations are the same feature.
\item Re-run Reticulum~II with equalisation enabled, for consistency with Section~\ref{sec:bootes}
      even though no contaminant is expected (Section~\ref{sec:config}).
\item Run the pooled cross-check on the Reticulum~II field (Section~\ref{sec:stackpool}).
\item Establish what distinguishes the fields in which the method recovers a known dwarf from those
      in which it recovers nothing. The batch outcome was strikingly bimodal and the cause is not
      known.
\end{itemize}

% =====================================================================
\appendix
\section{Hyperparameters}
\label{app:hyper}

Consolidated from Section~\ref{sec:config}, with justification.

\textbf{$h$, tile half-width.} Sets $n_Y$, and through the locality bound sets the achievable $K_M$
and hence the achievable significance (Sections~\ref{sec:km} and~\ref{sec:scan}). Smaller tiles are
not simply ``more local'': they are also less sensitive.

\textbf{$K_M$.} Bounded above by $0.05\min(n_X, n_Y)$ for locality and below by 25, since the
usable rank range is $20 \le K < K_M$. Larger is broadly better for sensitivity because $\Ups^*$
grows logarithmically while $\Ups_{\max}$ grows linearly; the cost is gradient leakage, which in
four dimensions is mild (Section~\ref{sec:km}). Must be shared across all eight references
(Section~\ref{sec:perref}), and resolved after equalisation.

\textbf{$p_{\rm ext}$.} A per-point exceedance level, not a cluster or field error rate
(Section~\ref{sec:upsstar}). Smaller tiles need a looser value: they carry a smaller
multiple-comparison burden, so the same physical excess sits at a less extreme quantile, and
simultaneously a lower $\Ups_{\max}$ ceiling. A diffuse overdensity is the case that most needs the
looser threshold, since it produces many mildly elevated points rather than a few extreme ones.

\textbf{$n_{\rm boot}$.} Controls how well the tail of the $\Gam$ null is resolved. At
$n_{\rm boot} = 10$ the quantile at $p_{\rm ext} = 10^{-3}$ rests on 20--30 null values
(Section~\ref{sec:gstar}).

\textbf{DBSCAN $\epsilon$, min.\ samples.} In scaled units, so their physical meaning depends on the
metric being frozen (Section~\ref{sec:scaling}).

\textbf{DPA $Z$.} Cluster-splitting parameter for the rep\^echage stage, left at the EagleEye
default.

\section{Parallelisation and determinism}
\label{app:parallel}

Three levels of parallelism exist and their product must stay within the core count:
configurations $\times$ references $\times$ inner $k$NN threads.

Only one backend works. Forking breaks OpenBLAS, because numpy initialises its thread pool in the
parent and the child inherits descriptors for threads that do not exist. Spawning re-imports
\texttt{\_\_main\_\_}, which is unusable from a notebook. Loky with the inner thread limit pinned to
one works, and is what is used. Note also that joblib disables nested parallelism, so requesting
eight configurations \emph{and} eight references does not give sixty-four workers --- the inner
level silently degrades.

Determinism required explicit care. Reseeding each bootstrap worker independently is not equivalent
to the serial computation: it moved $\Gstar$ from 52.38 to 61.66 in one test and changed the cluster
count from two to one. Advancing a single PCG64 stream to the offset each realisation would have
reached serially reproduces the serial result exactly. Measured speedups are 1.67$\times$ at
$n_{\rm boot} = 4$ and 3.47$\times$ at $n_{\rm boot} = 20$, with
$\max_i |\Delta\Gam_i| = 1.4\times10^{-14}$ against serial.

\section{Reserved}
\label{app:reserved}

Reserved for the estimator implementation item of Section~\ref{sec:eq9}, pending discussion with
the EagleEye authors.

\section{Figure manifest}
\label{app:figures}

All figures are symlinked into \texttt{report/figures/} from the run directories under
\texttt{field\_diagnostics/}, so the report is self-contained while the underlying products stay
with their runs.

\begin{table}[h]
\centering\small
\begin{tabular}{@{}lll@{}}
\toprule
File & Section & Source run\\
\midrule
\texttt{fig1\_reticulumII\_sky.png} & \ref{sec:retii} &
  Ret~II, $h$=1.3333, $K_M$=100, $p_{\rm ext}$=5e-4\\
\texttt{fig1b\_reticulumII\_pm.png} & \ref{sec:retii} & as above\\
\texttt{fig2\_bootesII\_sky.png} & \ref{sec:bootes}, \ref{sec:equalise} &
  Bo\"otes~II, $h$=1.0, $K_M$=50, $p_{\rm ext}$=1e-3, eq.\ on\\
\texttt{fig2b\_bootesII\_cmd.png} & \ref{sec:bootes} & as above\\
\texttt{fig3\_scan\_overlay.png} & \ref{sec:scan} &
  Ret~II scan, gid 0, $p_{\rm ext}$=5e-3, $K_M$ auto\\
\texttt{fig4\_reticulumII\_gamma\_null.png} & \ref{sec:gstar} & Ret~II, as fig.\ 1\\
\bottomrule
\end{tabular}
\end{table}

Figure 4 (the empirical $\Gam$ null against the analytic Fisher curves) is available and referenced
by Section~\ref{sec:gstar} but is not yet discussed in the text; adding that discussion is a small
outstanding item.

% =====================================================================
\section*{References}
\addcontentsline{toc}{section}{References}

\begin{list}{}{\leftmargin=1.6em \itemindent=-1.6em \itemsep=0.15em \parsep=0pt}
\item Battaglia G. et al., 2022, A\&A, 657, A54
\item Darragh-Ford E., Nadler E.~O., McLaughlin S., Wechsler R.~H., 2021, ApJ, 915, 48
\item Drlica-Wagner A. et al., 2020, ApJ, 893, 47
\item d'Errico M., Facco E., Laio A., Rodriguez A., 2021, Information Sciences, 560, 476
\item Ester M., Kriegel H.-P., Sander J., Xu X., 1996, Proc.\ KDD-96, 226
\item Fisher R.~A., 1932, Statistical Methods for Research Workers, 4th edn. Oliver \& Boyd,
      Edinburgh
\item Gaia Collaboration, 2023, A\&A, 674, A1
\item Glielmo A. et al., 2022, Patterns, 3, 100589
\item Guth A.~H., Namjoo M.~H., 2026, arXiv:2602.10178
\item Overdeck K. et al., 2026, arXiv:2606.09975
\item Rusterucci S., Hammer F., Yang Y., 2026, A\&A, accepted (arXiv:2608.21517)
\item Springer S., Scaffidi A., Autenrieth M., Contardo G., Laio A., Trotta R., Haario H., 2025,
      Scientific Reports (arXiv:2503.23927)
\item Tan C.~Y. et al., 2026, ApJ, 1000, 87
\item Vasiliev E., Baumgardt H., 2021, MNRAS, 505, 5978
\end{list}


\section{Chance clustering for a given test patch}

Throughout, the test patch is the central tile $\Yset$ of the $3\times3$ grid, and
all quantities are measured in physical (unscaled) units unless stated otherwise.

\begin{table}[h]
\centering
\caption{Notation. The symbol $\pmvec$ denotes proper motion throughout; the
Poisson mean is written $\nu$ rather than the more usual $\mu$ to avoid a clash.
Quantities marked \emph{measured} are taken from the data and require no model.}
\label{tab:notation}
\small
\begin{tabular}{@{}llp{8.1cm}@{}}
\toprule
Symbol & Units & Meaning \\
\midrule
\multicolumn{3}{@{}l}{\itshape Patch and field} \\
$n$              & ---            & number of stars in the test patch $\Yset$ \\
$A$              & $\mathrm{deg}^2$     & angular area of the test patch \\
$d$              & kpc       & assumed heliocentric distance of the candidate \\
$\pmvec=(\pmra,\pmdec)$ & $\mathrm{mas\,yr}^{-1}$ & proper-motion vector \\
$(\Delta\alpha\cos\delta,\,\Delta\delta)$ & deg & sky offsets from the patch centre \\
\midrule
\multicolumn{3}{@{}l}{\itshape Clump and cell} \\
$m$              & ---            & number of clump members inside the joint cell \\
$\Delta\mu$      & $\mathrm{mas\,yr}^{-1}$& radius of the proper-motion cell \\
$\theta$         & deg       & angular radius of the sky cell \\
$r_{\mathrm{phys}}$ & pc     & physical radius, $r_{\mathrm{phys}}=\theta\,d$ (small angle) \\
$q$              & ---            & aperture fraction: the cell radii are the $q$-quantiles of the member distances from the clump centroid \\
$N_{\mathrm{ap}}$& ---            & number of apertures $q$ scanned \\
\midrule
\multicolumn{3}{@{}l}{\itshape Statistics} \\
$f_\mu(\Delta\mu)$ & ---          & \emph{measured}: fraction of \emph{field} stars whose proper motion lies within $\Delta\mu$ of the clump's mean proper motion \\
$\nu$            & ---            & expected number of \emph{background} stars in the joint (PM $\times$ sky) cell \\
$N_{\mathrm{cells}}$ & ---        & number of independent joint cells tiling the patch; the look-elsewhere trials factor, $N_{\mathrm{cells}}=n/\nu$ \\
$P_{\mathrm{loc}}$ & ---          & local Poisson upper-tail probability, $P(\mathrm{Pois}(\nu)\ge m)$ \\
$\Nchance$       & ---            & expected number of chance clumps in the patch, look-elsewhere corrected \\
$\theta_{\mathrm{ur}}$, $r_{\mathrm{ur}}$ & deg, pc & ``unremarkable'' radius: the cell radius at which $\nu=m$ \\
\midrule
\multicolumn{3}{@{}l}{\itshape Concentration statistic (Sec.~\ref{sec:conc})} \\
$k$              & ---            & number of \emph{field} stars inside the proper-motion cell, $k \simeq n f_\mu$ \\
$c$              & ---            & largest number of those $k$ stars inside any sky disc of radius $\theta$ \\
$\nu_k$          & ---            & single-disc expectation for those $k$ stars, $\nu_k = k\pi\theta^2/A$ \\
$z$              & ---            & standardised scan statistic, $z=(c-\nu_k)/\sqrt{\nu_k}$ \\
\midrule
\multicolumn{3}{@{}l}{\itshape Pipeline quantities referenced} \\
$\Upsilon_i^{(j)}$ & ---          & EagleEye anomaly score of point $i$ against reference $j$ \\
$\Gamma_i$       & ---            & stacked score, $\Gamma_i=\sum_{j=1}^{R}\Upsilon_i^{(j)}$ \\
$A_\Gamma$       & ---            & set of points with $\Gamma_i>\Gamma^\star$ \\
$r_h$            & pc        & half-light radius of a stellar system \\
\bottomrule
\end{tabular}
\end{table}

% =====================================================================
\section{The separable null}
% =====================================================================

The hypothesis under test is not ``is there an overdensity'' but the narrower,
physically specific question:

\begin{quote}
\itshape
Given that the Galactic field has a strongly concentrated proper-motion
distribution, how often will $m$ stars that happen to share a proper-motion cell
\emph{also} fall within a small patch of sky, with no dynamical association
whatsoever?
\end{quote}

The corresponding null is \emph{separable}: proper motions are drawn from the
field's own (highly non-uniform) distribution, while sky positions are independent
and uniform over the patch,
%
\begin{equation}
p(\Delta\alpha\cos\delta,\,\Delta\delta,\,\pmvec)
\;=\;
\underbrace{\frac{1}{A}}_{\text{uniform on the sky}}
\times
\underbrace{p(\pmvec)}_{\text{measured from the field}} .
\label{eq:separable}
\end{equation}
%
Separability is what makes this the right null for the question: it retains
\emph{all} of the real proper-motion clustering and removes \emph{only} the
position--velocity correlation that a bound system would produce.

For a joint cell of proper-motion radius $\Delta\mu$ and sky radius $\theta$, the
expected number of background stars satisfying both conditions is
%
\begin{equation}
\boxed{\;
\nu(\Delta\mu,\theta)
\;=\;
n \;\cdot\; f_\mu(\Delta\mu) \;\cdot\; \frac{\pi\theta^{2}}{A}
\;}
\label{eq:nu}
\end{equation}
%
where the three factors are, respectively, the patch richness, the proper-motion
concentration, and the sky filling fraction. Crucially,
%
\begin{equation}
f_\mu(\Delta\mu)
\;=\;
\frac{1}{n_{\mathrm{bg}}}
\sum_{i \,\in\, \text{field}}
\mathbf{1}\!\left[\,\left\|\pmvec_i-\bar{\pmvec}_{\mathrm{clump}}\right\| \le \Delta\mu\,\right]
\label{eq:fmu}
\end{equation}
%
is \emph{measured}, with clump members excluded from the sum so that
$f_\mu$ estimates the background and not the signal. No Besan\c{c}on model, no
velocity ellipsoid, and no distance information enters.

The look-elsewhere-corrected expected number of chance clumps is then
%
\begin{equation}
\Nchance(\Delta\mu,\theta)
\;=\;
N_{\mathrm{cells}}\;P\!\left(\mathrm{Pois}(\nu) \ge m\right),
\qquad
N_{\mathrm{cells}} \;=\; \frac{A}{\pi\theta^{2}}\cdot\frac{1}{f_\mu} \;=\; \frac{n}{\nu} ,
\label{eq:Nchance}
\end{equation}
%
and a value $\Nchance \gtrsim 1$ means a clump this good is \emph{expected} in
this patch by chance.

\paragraph{Derivation of Eq.~\eqref{eq:Nchance}.}
Two independent routes agree. \emph{(i) Cell tiling.} Partition the occupied
product space into cells of the given size. The sky tiles $A/\pi\theta^{2}$ times;
a proper-motion cell holding a fraction $f_\mu$ of the stars tiles the occupied
proper-motion space $1/f_\mu$ times. Their product is $N_{\mathrm{cells}}=n/\nu$
by Eq.~\eqref{eq:nu}, and each cell independently has $\mathrm{Pois}(\nu)$
occupancy. \emph{(ii) Star-centred.} Place a cell on each star $i$; the remaining
$n-1$ stars enter independently with probability $\pi_C=\nu/n$, so the cell holds
$1+\mathrm{Bin}(n-1,\pi_C)$ stars and is flagged with probability
$P(\mathrm{Pois}(\nu)\ge m-1)$ --- note $m-1$, because the centre star is itself
one of the $m$. Linearity of expectation gives $n\,P(\mathrm{Pois}(\nu)\ge m-1)$
flagged centres \emph{exactly}, despite the cells overlapping; a genuine clump is
flagged from each of its $m$ members, so dividing by $m$ de-duplicates. At the
best aperture of Table~\ref{tab:aperture} route (i) gives $0.155$ and route (ii)
$0.160$, agreeing to $3\%$. Pairing route (ii)'s trials factor $n/m$ with route
(i)'s tail $P(\ge m)$ --- a tempting but inconsistent hybrid --- discards a factor
$\nu/m$ and understates $\Nchance$ by $2.8\times$ here.

\paragraph{Aperture scan.}
Reporting the minimum over $N_{\mathrm{ap}}$ apertures incurs a trials penalty
bounded above by Bonferroni. The apertures are \emph{nested}, hence strongly
positively correlated, so the effective number of independent trials is
considerably smaller than $N_{\mathrm{ap}}$ --- of order $2$--$3$ for the seven
used here. We therefore quote the bracket
%
\begin{equation}
\Nchance \;\le\; \Nchance^{\,\mathrm{scan}} \;\le\; N_{\mathrm{ap}}\,\Nchance .
\label{eq:scanbracket}
\end{equation}

\paragraph{Product cell versus joint ball.}
Equation~\eqref{eq:nu} defines a \emph{product} cell --- ``same proper motion
\textbf{and} same patch of sky'' --- with a fixed, interpretable shape. This is
deliberately distinct from a ball in the MAD-scaled four-dimensional metric used
internally by EagleEye. A joint ball \emph{couples} the position and
proper-motion axes, and a nearest-neighbour construction within it selects
whichever subset of stars is most concentrated in that coupled metric. For the
same stars the two constructions can differ by many orders of magnitude; the
product cell is both the more conservative and the more physically interpretable,
and is the one quoted here.

\paragraph{What this statistic does \emph{not} measure.}
Equation~\eqref{eq:Nchance} is a \emph{counting} statistic: it asks only whether
at least $m$ stars fall inside the cell, and is blind to how they are arranged
within it. A tight knot at the cell centre and a uniform filling of the cell
receive identical scores. Concentration enters only indirectly, through the cell
radii being set by the observed member distances and then scanned --- a crude
profile fit. Section~\ref{sec:conc} supplements it with a statistic that measures
concentration directly.

% =====================================================================
\section{Application to the Reticulum~II field}
% =====================================================================

Patch parameters: $n=1595$, $A=2.29\degsq$, assumed $d=30\unit{kpc}$
(so $1^\circ=524\unit{pc}$). The candidate is the single $\Gamma$ cluster of
the run, with 39 members. Cell radii are set by the $q$-quantiles of the member
distances from the clump centroid, separately in the proper-motion and sky
subspaces.

\begin{table}[h]
\centering
\caption{Chance-clump expectation as a function of aperture. Each row uses cell
radii $(\Delta\mu,\theta)$ equal to the $q$-quantiles of the member distances,
retaining the $m$ members inside \emph{both}. $f_\mu$ is measured with the 39
clump members excluded from the background. The last two columns bracket the
aperture-scan penalty as in Eq.~\eqref{eq:scanbracket}. The minimum, marked
$\dagger$, is the most favourable reading available and still gives $\Nchance$
between $0.15$ and $1.1$.}
\label{tab:aperture}
\small
\begin{tabular}{@{}ccccc rrrrrr@{}}
\toprule
$q$ & $m$ & $\Delta\mu$ & $\theta$ & $r_{\mathrm{phys}}$ & $f_\mu$ & $\nu$
& $N_{\mathrm{cells}}$ & $P_{\mathrm{loc}}$ & $\Nchance$ & $N_{\mathrm{ap}}\Nchance$ \\
    &     & [$\mathrm{mas\,yr}^{-1}$] & [deg] & [pc] & & & & & & \\
\midrule
0.4 &  8 & 0.409 & 0.194 & 101 & 0.031 &   2.52 & 632 & $4.48\times10^{-3}$ & 2.83 & 19.8 \\
0.5 & 12 & 0.475 & 0.214 & 112 & 0.040 &   3.95 & 404 & $8.16\times10^{-4}$ & 0.330 & 2.31 \\
$0.6^{\dagger}$ & 15 & 0.542 & 0.231 & 121 & 0.048 & 5.43 & 294 & $5.27\times10^{-4}$ & \textbf{0.155} & \textbf{1.08} \\
0.7 & 19 & 0.604 & 0.270 & 142 & 0.055 &   8.63 & 185 & $1.54\times10^{-3}$ & 0.284 & 1.99 \\
0.8 & 27 & 0.734 & 0.318 & 167 & 0.071 &  15.47 & 103 & $4.91\times10^{-3}$ & 0.506 & 3.54 \\
0.9 & 33 & 0.948 & 0.398 & 208 & 0.095 &  32.20 &  50 & $4.68\times10^{-1}$ & 23.2 & 162 \\
1.0 & 39 & 1.882 & 0.549 & 287 & 0.213 & 136.90 &  12 & $1.00\times10^{0}$  & 11.7 & 81.6 \\
\bottomrule
\end{tabular}
\end{table}

\paragraph{The proper-motion enhancement.}
At the best aperture the clump's proper-motion cell is a disc of radius
$0.542\masyr$, covering $\pi(0.542)^2/100 = 0.92\%$ of the
$\pm5\masyr$ proper-motion box, yet it contains $f_\mu=4.76\%$ of the
field. The field is therefore
%
\begin{equation}
\frac{f_\mu}{\pi\Delta\mu^{2}/A_{\pmvec}} \;=\; \frac{0.0476}{0.0092} \;\simeq\; 5.2
\end{equation}
%
times denser in that proper-motion cell than a uniform distribution would give.
This is precisely the effect of concern: \emph{proper-motion concentration makes
sky clumping about five times cheaper than a naive uniform estimate}, and it is
quantified here without recourse to any Galactic model.

% =====================================================================
\section{A statistic that measures concentration}
\label{sec:conc}
% =====================================================================

Take the $k$ stars within $\Delta\mu$ of the clump's mean proper motion, and let
$c$ be the largest number of them falling inside any disc of radius $\theta$ on
the sky. Standardising by the single-disc expectation $\nu_k=k\pi\theta^{2}/A$
gives
%
\begin{equation}
z \;=\; \frac{c-\nu_k}{\sqrt{\nu_k}} .
\label{eq:zscan}
\end{equation}
%
This responds to \emph{where} the stars sit, not merely how many are present, and
is standardised so that proper-motion cells of differing occupancy may be
compared. Critically, $z$ must be calibrated empirically: because $c$ is a
\emph{maximum} over disc placements, $\mathbb{E}[z]>0$ even for perfectly uniform
points, and the offset is large.

\begin{table}[h]
\centering
\caption{Sky concentration of proper-motion-selected subsamples, measured by
Eq.~\eqref{eq:zscan} at $\Delta\mu=0.542\,\mathrm{mas\,yr}^{-1}$,
$\theta=0.231^\circ$. The three nulls are: other proper-motion cells of the same
radius drawn from the field; $k=97$ points thrown uniformly on the patch; and the
real sky positions with proper motions permuted. All three give a median of
$z\simeq2.1$--$2.3$, which is therefore the look-elsewhere offset intrinsic to a
maximum over disc placements and \emph{not} evidence of real clustering. The
observed $z=2.95$ is unremarkable against every null.}
\label{tab:conc}
\small
\begin{tabular}{@{}lrrrrr@{}}
\toprule
Sample & trials & median $z$ & 90\% & max $z$ & $p(\ge z_{\mathrm{obs}})$ \\
\midrule
\textbf{Observed clump} & --- & \multicolumn{4}{l}{$z_{\mathrm{obs}}=2.95$ \quad($c=15$, $k=97$, $\nu_k=7.12$)} \\
\midrule
Other PM cells, field       & 300 & 2.14 & 2.97 & 3.85 & 0.116 \\
Uniform points, $k=97$      & 300 & 2.20 & 3.33 & 4.45 & 0.223 \\
PM shuffled, real sky       & 172 & 2.30 & 3.25 & 5.53 & 0.214 \\
\midrule
Poisson prediction          & --- & 0    & ---  & ---  & $6.6\times10^{-3}$ \\
\bottomrule
\end{tabular}
\end{table}

Two conclusions follow. First, the observed concentration is reached by
$12$--$22\%$ of randomly placed proper-motion cells, so the clump is not unusually
concentrated. Second, the single-disc Poisson probability ($6.6\times10^{-3}$) is
optimistic by a factor of $\sim20$ relative to the empirical rate; the discrepancy
is the look-elsewhere effect of maximising over disc placements, and is precisely
what the $N_{\mathrm{cells}}$ factor of Eq.~\eqref{eq:Nchance} exists to absorb.
Indeed the implied effective trials, $0.15/6.6\times10^{-3}\approx23$, agree in
order of magnitude with the sky-only tiling number $A/\pi\theta^{2}\approx14$. The
uniform control is indispensable: without it the median offset of $z\simeq2.2$
would be mistaken for intrinsic clumpiness of the Galactic field.

% =====================================================================
\section{Dependence on patch richness}
% =====================================================================

Holding the clump fixed at the best aperture ($m=15$, $\theta=0.231^\circ$,
$f_\mu=0.048$) and varying only the patch richness $n$ separates two effects that
are easily conflated.

\begin{table}[h]
\centering
\caption{Chance-clump expectation versus patch richness, for a fixed clump
($m=15$, $\theta=0.231^\circ$, $f_\mu=0.048$, $A=2.29\degsq$,
$d=30\unit{kpc}$). Columns 2--4 hold the \emph{angular} scale fixed; columns 5--6
instead hold the \emph{member count} fixed and solve $\nu=m$ for the radius at
which the clump ceases to be remarkable. The two halves answer different
questions and point in opposite directions.}
\label{tab:richness}
\small
\begin{tabular}{@{}r rrr rr@{}}
\toprule
$n$ & $\nu$ & $P_{\mathrm{loc}}$ & $\Nchance$
& $\theta_{\mathrm{ur}}$ [deg] & $r_{\mathrm{ur}}$ [pc] \\
\midrule
 200 &  0.70 & $1.81\times10^{-15}$ & $2.42\times10^{-14}$ & 1.071 & 560.9 \\
 500 &  1.75 & $6.35\times10^{-10}$ & $2.12\times10^{-8}$  & 0.678 & 354.7 \\
1000 &  3.49 & $4.13\times10^{-6}$  & $2.75\times10^{-4}$  & 0.479 & 250.8 \\
2000 &  6.98 & $5.58\times10^{-3}$  & $7.44\times10^{-1}$  & 0.339 & 177.4 \\
4000 & 13.96 & $4.26\times10^{-1}$  & $1.14\times10^{2}$   & 0.240 & 125.4 \\
8000 & 27.92 & $9.97\times10^{-1}$  & $5.32\times10^{2}$   & 0.169 &  88.7 \\
\bottomrule
\end{tabular}
\end{table}

At \emph{fixed angular scale} (column 4), sparse patches produce \emph{fewer}
chance clumps: there are simply fewer background stars available to fake one, and
$\Nchance$ falls by fourteen orders of magnitude between $n=8000$ and $n=200$.

At \emph{fixed member count}, however, the radius at which a clump stops being
remarkable grows as
%
\begin{equation}
\theta_{\mathrm{ur}} \;=\; \sqrt{\frac{m A}{\pi n f_\mu}} \;\propto\; n^{-1/2} .
\label{eq:thetaur}
\end{equation}
%
Because EagleEye's $k$-nearest-neighbour balls are adaptive, a sparse patch drives
the method to physically enormous apertures. At $n=200$ any 15-star grouping out
to $r_{\mathrm{ur}}=561\unit{pc}$ is consistent with background. Since ultra-faint
dwarfs have $r_h \sim 20\text{--}300\unit{pc}$ (Reticulum~II itself has
$r_h\approx50\unit{pc}$), the method in that regime cannot distinguish a bound
system from the field \emph{at any reported significance}, because the statistic
has no knowledge of parsecs.

\medskip
\noindent
This is the useful physical bound: compare $r_{\mathrm{ur}}$ from
Eq.~\eqref{eq:thetaur} against $r_h$ for the systems of interest. Where the two
overlap, astrometric significance alone is uninformative.

% =====================================================================
\section{Summary}
% =====================================================================

\begin{enumerate}\itemsep2pt
\item Under the separable null of Eq.~\eqref{eq:separable}, the observed clump has
      $0.15 \le \Nchance^{\,\mathrm{scan}} \le 1.1$ (Table~\ref{tab:aperture}):
      of order unity, and therefore not on its own evidence of a bound system.
\item The counting statistic of Eq.~\eqref{eq:Nchance} is blind to concentration.
      A direct concentration statistic (Eq.~\ref{eq:zscan}) gives
      $p=0.12$--$0.22$ against three independent nulls (Table~\ref{tab:conc}),
      confirming the conclusion by an independent route.
\item The proper-motion effect is real and quantifiable: the clump's proper-motion
      cell is $5.2\times$ overdense relative to a uniform proper-motion
      distribution, so sky clumping is correspondingly cheap.
\item Sparse patches do not produce more chance clumps; they produce chance clumps
      that are \emph{too large to be bound systems} (Table~\ref{tab:richness}).
\item The separable null assumes position--proper-motion independence. If that is
      violated in the patch, the true chance rate is \emph{higher}, not lower.
      The eight tiled reference sets $\Xset^{(j)}$ provide the data to test it.
\end{enumerate}

\paragraph{Implementation.}
The three estimates are produced by
\begin{itemize}\itemsep0pt
\item \texttt{wake\_test.pm\_sky\_chance\_clumps}\quad(Table~\ref{tab:aperture}),
\item \texttt{wake\_test.pm\_sky\_concentration\_test}\quad(Table~\ref{tab:conc}),
\item \texttt{wake\_test.chance\_vs\_patch\_richness}\quad(Table~\ref{tab:richness}).
\end{itemize}



\end{document}